\section{Marco teórico}
\label{sec:marco_teorico}

\subsection{Notación}

Sea una pregunta con $n$ opciones etiquetadas ($n$ es la cardinalidad; el primitivo \texttt{choice} admite
hasta 255 y \texttt{score} entre 2 y 10). El modelo devuelve una distribución $p = (p_1, \dots, p_n)$ sobre
las opciones, con $p_i \ge 0$ y $\sum_i p_i = 1$, y la etiqueta correcta es $y \in \{1,\dots,n\}$.
Escribimos $p_{\max} = \max_i p_i$, $\hat{y} = \arg\max_i p_i$, y el acierto $a = \mathbf{1}[\hat{y} = y]$.
Con $H(p) = -\sum_i p_i \log p_i$ denotamos la entropía de Shannon (el logaritmo es indistinto mientras se
normalice). Ordenamos las probabilidades de mayor a menor: $p_{(1)} \ge p_{(2)} \ge \dots \ge p_{(n)}$, de
modo que $p_{(1)} = p_{\max}$.

\subsection{El campo \texttt{confidence} como función de la distribución}

Cada implementación define \texttt{confidence} como una función distinta de $p$ y de $n$ (Tabla~\ref{tab:campo_confidence}):

\begin{table}[t]
\centering
\caption{Definiciones de \texttt{confidence} como función de la distribución $p$ y de la cardinalidad $n$.}
\label{tab:campo_confidence}
\footnotesize
\resizebox{\columnwidth}{!}{%
\begin{tabular}{@{}llll@{}}
\toprule
definición & fórmula & dominio & quién la usa \\
\midrule
concentración de Jev & $C_{\mathrm{Jev}}(p) = (n\,p_{\max} - 1)/(n - 1)$ & $[0,1]$ & Jev, Kev \\
entropía normalizada & $C_{\mathrm{ent}}(p) = 1 - H(p)/\log n$ & $[0,1]$ & Laya (\texttt{choice}, \texttt{score}) \\
probabilidad ganadora & $C_{\max}(p) = p_{\max}$ & $[1/n, 1]$ & jev-lite, \texttt{answer\_confidence} \\
margen & $C_{\mathrm{mar}}(p) = p_{(1)} - p_{(2)}$ & $[0,1]$ & derivados \\
\texttt{noul} (binario) & $C_{\mathrm{noul}}(p) = \max(p, 1 - p)$ & $[1/2, 1]$ & Laya (\texttt{noul}) \\
\bottomrule
\end{tabular}}
\end{table}

Para una misma $p$, estas funciones toman valores distintos. Esta es la fuente de la incompatibilidad de
umbrales: el umbral es un punto en el rango de una función, y los rangos no coinciden (Tabla~\ref{tab:cuatro_numeros}).

\begin{table}[t]
\centering
\caption{La misma distribución, cuatro definiciones, cuatro números.}
\label{tab:cuatro_numeros}
\small
\begin{tabular}{lcccc}
\toprule
$p$ ($n = 3$) & $C_{\mathrm{Jev}}$ & $C_{\mathrm{ent}}$ & $C_{\max}$ & $C_{\mathrm{mar}}$ \\
\midrule
$(0.0118, 0.4629, 0.5253)$ & 0.288 & 0.320 & 0.525 & 0.062 \\
\bottomrule
\end{tabular}
\end{table}

\subsection{Calibración frente a discriminación}

La calibración pregunta si el número es \emph{creíble}: $P(\hat{y} = y \mid C = c) = c$ para todo
$c$. La discriminación pregunta si el número \emph{ordena bien}: si los aciertos reciben valores de
$C$ más altos que los errores, y se mide con AUROC$(C, a)$. Un modelo puede discriminar bien y estar
des-calibrado, y dos definiciones pueden discriminar idéntico y calibrar distinto: por eso el trabajo
reporta ambas. El ECE (Error de Calibración Esperado) con $M$ bins $B_m$ es

\[
\mathrm{ECE} = \sum_m \frac{\lvert B_m\rvert}{N}\cdot \lvert \mathrm{acc}(B_m) - \mathrm{conf}(B_m)\rvert,
\]

y se calcula en dos esquemas, bins de igual ancho y de igual masa, porque el número
cambia entre ellos \cite{naeini2015,guo2017}. Como métricas de error probabilístico se usan la
Brier multiclase $\mathrm{Brier} = (1/N)\sum_i \sum_k (p_{ik} - y_{ik})^2$ \cite{brier1950} y la
NLL $-(1/N)\sum_i \log p_i(y_i)$.

\subsection{Equivalencia en orden y transferencia exacta}

El punto teórico del trabajo es que la incompatibilidad no es uniforme: algunas definiciones están
relacionadas por una transformación monótona, y entonces el umbral transfiere exacto; otras no lo
están. Lo fijamos en tres proposiciones (demostraciones y comprobación numérica en
\texttt{paper/03\_especificacion\_confidence.md}).

\begin{quote}
\textbf{Proposición 1 (equivalencia en orden).} Sea $n$ fijo. La definición de Jev es afín y estrictamente
creciente en $p_{\max}$: $C_{\mathrm{Jev}} = \alpha\,p_{\max} + \beta$ con $\alpha = n/(n-1) > 0$ y
$\beta = -1/(n-1)$. En consecuencia, $C_{\mathrm{Jev}}$ y $C_{\max}$ son \textbf{equivalentes en orden}:
existe $\phi$ estrictamente creciente con $C_{\mathrm{Jev}} = \phi(C_{\max})$. Un umbral $\tau$ sobre
$C_{\mathrm{Jev}}$ induce el mismo conjunto de actuaciones y abstenciones que el umbral $\phi^{-1}(\tau)$
sobre $C_{\max}$, y la cobertura es idéntica.
\end{quote}

\begin{quote}
\textbf{Proposición 2 (no equivalencia con la entropía).} Para $n \ge 3$ existen distribuciones $p, p'$ con
$p_{\max} = p'_{\max}$ y $H(p) \ne H(p')$. Por tanto $C_{\mathrm{ent}}$ \textbf{no} es función de
$p_{\max}$ solo, y en general \textbf{no} existe $\phi$ monótona con $C_{\mathrm{ent}} = \phi(C_{\max})$. El
umbral no transfiere entre una definición tipo $p_{\max}$ y una tipo entropía.
\end{quote}

\begin{quote}
\textbf{Proposición 3 (degeneración en $n = 2$).} Con $n = 2$, $C_{\mathrm{Jev}} = 2\,p_{\max} - 1 =
C_{\mathrm{mar}}$, esto es, la definición de Jev y el margen son la \textbf{misma} función; y
$C_{\mathrm{ent}}$, con $H$ función de un solo parámetro, es también una transformación estrictamente
monótona de $p_{\max}$. Por tanto en $n = 2$ \textbf{todas} las definiciones son equivalentes en orden, y la
incompatibilidad de umbrales es un fenómeno exclusivo de $n \ge 3$ con distribuciones dispersas.
\end{quote}

\textbf{Consecuencia: el fallo tiene dos capas.} Si dos definiciones son equivalentes en orden pero viven en
escalas distintas (por ejemplo $C_{\mathrm{Jev}}$ y $C_{\max}$), el número no viaja aunque el
mapa sí: existe traducción exacta $\phi^{-1}$. Si además pertenecen a clases distintas
(por ejemplo $p_{\max}$ y la entropía), no existe ni número ni mapa, y el porte directo es una extrapolación
inválida. Esta distinción, entre la capa de escala y la capa de forma, es la que explica
por qué el resultado no se puede resumir con un promedio y se documenta en
\texttt{docs/erratas\_especificacion.md}.

\subsection{Transferencia de umbrales (definición formal)}

Dados un origen $A$, un destino $B$, una cobertura objetivo $q \in (0, 1)$ y un conjunto de calibración:

\begin{enumerate}
\item Hallar $\tau_A$ tal que $\mathrm{cobertura}_A(\tau_A) = q$, con
  $\mathrm{cobertura}_A(\tau) = \#\{C_A \ge \tau\}/N$ (el menor umbral observado que no supera la cobertura
  objetivo).
\item Portarlo al destino de una de dos maneras:
  \begin{itemize}
  \item porte directo: $\tau_B := \tau_A$, que es lo que practica el ecosistema;
  \item porte con mapa: $\tau_B := \psi(\tau_A)$, con $\psi$ el mapa cerrado $A \to B$, que
    solo existe dentro de una clase de equivalencia y a cardinalidad fija (por ejemplo
    $\psi = \phi^{-1}$ para \texttt{jev} $\to$ \texttt{pmax}).
  \end{itemize}
\item Medir en el conjunto de prueba la cobertura efectiva $\mathrm{cobertura}_B(\tau_A)$ y la precisión.
\end{enumerate}

La degradación $\Delta(A \to B) = q - \mathrm{cobertura}_B(\tau_A)$ es la celda de la matriz. Se
repite por modelo, por dataset y por cardinalidad. El mismo procedimiento se puede aplicar entre
datasets (umbral ajustado en un dataset $A$ y aplicado en otro $B$), que es el eje donde \cite{jebadiah}
observó que la calibración no viaja y \cite{arXiv2609.33843} lo cuantificó dentro de la definición
$p_{\max}$. Este trabajo recorre el eje distinto, el de las definiciones.

\subsection{Predicción selectiva: el umbral como política}

Con probabilidad de acierto $p$, coste de error $C_{\mathrm{err}}$ y coste de escalar $C_{\mathrm{esc}}$,
actuar conviene si $(1-p)\cdot C_{\mathrm{err}} < C_{\mathrm{esc}}$, es decir
$p > 1 - C_{\mathrm{esc}}/C_{\mathrm{err}}$ \cite{chow1970}. Esta frontera presupone que $p$ es una
probabilidad calibrada. Cuando el umbral se elige como cuantil de una \texttt{confidence} mal especificada,
la política efectiva deja de ser la óptima para el coste real: la métrica de la matriz es, en el fondo, la
pérdida de utilidad del despliegue.

\subsection{Nota sobre RLCD}

RLCD (ruido gaussiano de media cero en los logits, regla estrictamente propia, REINFORCE con baseline de
grupo, temperatura post-hoc) es el mecanismo que produce estas \texttt{confidence}. La teoría garantiza que
reportar honestamente es el óptimo bajo una regla propia \cite{gneiting2007}; el análisis de septiembre de
2026 muestra que, en la práctica, el término RL equivale al gradiente de una cross-entropy suavizada con
ruido, de modo que la inferencia sin ruido sale sobre-confiada
\cite{arXiv2609.38850,lakoreai2026}. Esto explica por qué hace falta medir la calibración, pero no
es el objeto de este trabajo: aquí no se propone un objetivo de entrenamiento, se construye el instrumento
que mide su resultado.
